\documentclass[aspectratio=169]{beamer}

% ================================================================
% TEMA Y COLORES
% ================================================================
\usetheme{Madrid}
\usecolortheme{default}

\usepackage[utf8]{inputenc}
\usepackage[T1]{fontenc}
\usepackage[spanish]{babel}
\usepackage{amsmath, amssymb}
\usepackage{booktabs}
\usepackage{array}
\usepackage{listings}
\usepackage{pgfplots}
\pgfplotsset{compat=1.18}
\usepackage{tikz}
\usetikzlibrary{arrows.meta, calc, patterns, decorations.pathreplacing}
\usepackage{float}

% ----------------------------------------------------------------
% Colores del documento (idénticos a memoria.tex)
% ----------------------------------------------------------------
\definecolor{myblue}{RGB}{94, 52, 235}
\definecolor{myyellow}{RGB}{211, 121, 32}
\definecolor{myred}{RGB}{222, 75, 49}
\definecolor{mygreen}{RGB}{23, 110, 63}

\definecolor{backcolour}{rgb}{0.96,0.96,0.96}
\definecolor{codegreen}{rgb}{0.13,0.55,0.13}
\definecolor{codegray}{rgb}{0.5,0.5,0.5}
\definecolor{codepurple}{rgb}{0.58,0,0.82}
\definecolor{codeblue}{rgb}{0.0,0.0,0.8}

% ----------------------------------------------------------------
% Paleta Beamer basada en myblue
% ----------------------------------------------------------------
\setbeamercolor{palette primary}  {bg=myblue,      fg=white}
\setbeamercolor{palette secondary}{bg=myblue!80,   fg=white}
\setbeamercolor{palette tertiary} {bg=myblue!60,   fg=white}
\setbeamercolor{palette quaternary}{bg=myblue!40,  fg=white}
\setbeamercolor{structure}        {fg=myblue}
\setbeamercolor{frametitle}       {bg=myblue, fg=white}
\setbeamercolor{title}            {bg=myblue, fg=white}
\setbeamercolor{block title}      {bg=myblue!85, fg=white}
\setbeamercolor{block body}       {bg=myblue!8}
\setbeamercolor{block title alerted}{bg=myred!85, fg=white}
\setbeamercolor{block body alerted} {bg=myred!8}
\setbeamercolor{block title example}{bg=mygreen!75, fg=white}
\setbeamercolor{block body example} {bg=mygreen!8}
\setbeamercolor{itemize item}     {fg=myblue}
\setbeamercolor{itemize subitem}  {fg=myyellow}

\setbeamerfont{frametitle}{series=\bfseries}
\setbeamerfont{title}{series=\bfseries}

% Quitar botones de navegación
\setbeamertemplate{navigation symbols}{}
% Número de página
\setbeamertemplate{footline}[frame number]

% ----------------------------------------------------------------
% lstlisting (pseudocódigo) — mismo estilo que memoria.tex
% ----------------------------------------------------------------
\lstdefinestyle{modern}{
    backgroundcolor=\color{backcolour},
    commentstyle=\color{codegreen}\itshape,
    keywordstyle=\color{codeblue}\bfseries,
    numberstyle=\tiny\color{codegray},
    stringstyle=\color{codepurple},
    basicstyle=\ttfamily\scriptsize,
    breakatwhitespace=false,
    breaklines=true,
    captionpos=b,
    keepspaces=true,
    numbers=left,
    numbersep=8pt,
    showspaces=false,
    showstringspaces=false,
    showtabs=false,
    tabsize=4,
    frame=single,
    rulecolor=\color{lightgray},
    xleftmargin=12pt,
    morekeywords={Mientras, hacer, Para, hasta, Si, entonces, Si_no,
                  Fin, Retornar, Inicializar, Evaluar, de},
    morecomment=[s]{(*}{*)},
    sensitive=false,
    literate={á}{{\'a}}1 {é}{{\'e}}1 {í}{{\'i}}1 {ó}{{\'o}}1 {ú}{{\'u}}1
             {Á}{{\'A}}1 {É}{{\'E}}1 {Í}{{\'I}}1 {Ó}{{\'O}}1 {Ú}{{\'U}}1
             {ñ}{{\~n}}1 {Ñ}{{\~N}}1
}
\lstset{style=modern}

% ----------------------------------------------------------------
% Comandos auxiliares
% ----------------------------------------------------------------
\newcommand{\highlight}[1]{\textcolor{myblue}{\textbf{#1}}}
\newcommand{\warn}[1]{\textcolor{myred}{\textbf{#1}}}
\newcommand{\good}[1]{\textcolor{mygreen}{\textbf{#1}}}
\renewcommand{\accent}[1]{\textcolor{myyellow}{\textbf{#1}}}

% ================================================================
% METADATOS
% ================================================================
\title[SCA]{Sine Cosine Algorithm (SCA)}
\subtitle{Una metaheurística basada en funciones trigonométricas}
\author{Jesús Muñoz Velasco}
\institute{Metaheurísticas · Trabajo Final}
\date{Curso 2025--2026}

% ================================================================
\begin{document}

% ----------------------------------------------------------------
% PORTADA
% ----------------------------------------------------------------
\begin{frame}[plain]
  \titlepage
\end{frame}

% ----------------------------------------------------------------
% ÍNDICE
% ----------------------------------------------------------------
\begin{frame}{Índice}
  \tableofcontents
\end{frame}

% ================================================================
\section{Introducción}
% ================================================================

\begin{frame}{¿Qué es el SCA?}
  \begin{columns}[T]
    \column{0.55\textwidth}
      \begin{block}{Origen}
        Propuesto por \textbf{Seyedali Mirjalili} en 2016\\
        \textit{Knowledge-Based Systems}, vol.~96, pp.~120--133
      \end{block}
      \vspace{0.4em}
      \begin{itemize}
        \item Metaheurística \highlight{poblacional} basada en modelos matemáticos (\textit{math-inspired})
        \item Aprovecha las propiedades geométricas de \highlight{seno} y \highlight{coseno}
        \item Sin inspiración biológica ni física: pura matemática
      \end{itemize}

    \column{0.42\textwidth}
      \begin{exampleblock}{Idea central}
        $\sin$ y $\cos$ oscilan periódicamente entre $[-1,1]$\\[0.4em]
        $\Rightarrow$ Si se \textbf{escalan} con $r_1 > 1$: exploración\\
        $\Rightarrow$ Si $r_1 < 1$: explotación
      \end{exampleblock}
      \vspace{0.4em}
      \begin{alertblock}{Ventaja clave}
        El balance exploración/explotación emerge \textbf{automáticamente} del decaimiento de $r_1$
      \end{alertblock}
  \end{columns}
\end{frame}

% ================================================================
\section{Fundamentos matemáticos}
% ================================================================

\begin{frame}{Ecuación de actualización}
  \begin{block}{Regla de movimiento (dimensión $j$, iteración $t$)}
    \[
      X_j^{t+1} =
      \begin{cases}
        X_j^t + r_1 \cdot \sin(r_2) \cdot \bigl|r_3 P_j^t - X_j^t\bigr|, & r_4 < 0{,}5 \\[6pt]
        X_j^t + r_1 \cdot \cos(r_2) \cdot \bigl|r_3 P_j^t - X_j^t\bigr|, & r_4 \geq 0{,}5
      \end{cases}
    \]
  \end{block}

  \vspace{0.5em}
  Factorizando con $\phi = r_1 \cdot f(r_2)$ y $|\Delta_j^t| = |r_3 P_j^t - X_j^t|$:
  \[
    \boxed{X_j^{t+1} = X_j^t + \phi \cdot |\Delta_j^t|}
  \]

  \vspace{0.3em}
  \begin{center}
    \small
    \begin{tabular}{c l l}
      \toprule
      \textbf{Param.} & \textbf{Rango} & \textbf{Rol} \\
      \midrule
      $r_1$ & $[0,\,2]$ decreciente & \highlight{Amplitud}: exploración vs.\ explotación \\
      $r_2$ & $[0,\,2\pi]$ aleatorio & Dirección y signo del paso \\
      $r_3$ & $[0,\,2]$ aleatorio & Peso estocástico de la distancia al destino \\
      $r_4$ & $[0,\,1]$ aleatorio & Selección equiprobable sin/cos \\
      \bottomrule
    \end{tabular}
  \end{center}
\end{frame}

% ================================================================
\section{Exploración vs.\ explotación}
% ================================================================

% ----------------------------------------------------------------
% 3.1 Mapa de componentes
% ----------------------------------------------------------------
\begin{frame}{¿Qué componente provoca qué?}
  \begin{columns}[T]
    \column{0.48\textwidth}
      \begin{block}{\good{Componentes de exploración}}
        \begin{itemize}
          \item \textbf{$r_1 > 1$}: amplifica el paso más allá del destino
          \item \textbf{$r_3 > 1$}: aleja el destino efectivo, agranda $|\Delta_j^t|$
          \item \textbf{$r_2$ uniforme en $[0,2\pi]$}: cubre todas las direcciones, incluidas las de signo negativo ($f(r_2) < 0$)
          \item \textbf{$r_4$ aleatorio}: seno y coseno desfasados $\pi/2$, mejoran la isotropía
        \end{itemize}
      \end{block}

    \column{0.48\textwidth}
      \begin{alertblock}{\accent{Componentes de explotación}}
        \begin{itemize}
          \item \textbf{$r_1 < 1$}: confina el paso al interior de $[X_j^t,\, P_j^t]$
          \item \textbf{$r_3 < 1$}: acerca el destino efectivo a $X_j^t$, reduce $|\Delta_j^t|$
          \item \textbf{Punto destino global $P_j^t$}: atractor común que concentra la población
          \item \textbf{Decaimiento de $r_1$}: garantiza convergencia progresiva
        \end{itemize}
      \end{alertblock}
  \end{columns}

  \vspace{0.5em}
  \begin{center}
    \small
    $r_1$ es el \highlight{único parámetro determinista}; el resto son aleatorios $\Rightarrow$ el balance es \textbf{estocástico pero dirigido}
  \end{center}
\end{frame}

% ----------------------------------------------------------------
% 3.2 Figura r1
% ----------------------------------------------------------------
\begin{frame}{$r_1$: el director del equilibrio}
  \begin{columns}[T]
    \column{0.58\textwidth}
      \begin{tikzpicture}
      \begin{axis}[
        width=\textwidth, height=5.8cm,
        axis lines*=left,
        axis line style={draw=black, line width=0.6pt},
        xlabel={$t$}, ylabel={$r_1$},
        label style={font=\small\sffamily},
        tick label style={font=\scriptsize\sffamily},
        xmin=0, xmax=100, ymin=0, ymax=2.15,
        xtick={0,25,50,75,100}, ytick={0,0.5,1,1.5,2},
        tick align=outside,
        grid=both,
        grid style={line width=0.4pt, draw=gray!30},
        legend style={font=\scriptsize\sffamily, at={(0.97,0.97)},
          anchor=north east, draw=none, fill=white,
          fill opacity=0.85, text opacity=1, row sep=2pt},
        legend cell align={left},
      ]
        \addplot[very thick, myblue, domain=0:100, samples=200]{2*(1 - x/100)};
        \addlegendentry{$r_1(t)=2(1-t/T)$}
        \addplot[very thick, dashed, myred, domain=0:100, samples=2]{1};
        \addlegendentry{Umbral $r_1=1$}
        \addplot[fill=mygreen!40, opacity=0.6, draw=none]
          coordinates {(0,2)(50,1)(50,1)(0,1)} \closedcycle;
        \addplot[fill=myyellow!40, opacity=0.6, draw=none]
          coordinates {(50,1)(100,0)(100,1)(50,1)} \closedcycle;
        \node[font=\small\bfseries, mygreen]  at (axis cs:18, 1.55) {Exploración};
        \node[font=\small\bfseries, myyellow] at (axis cs:78, 0.38) {Explotación};
      \end{axis}
      \end{tikzpicture}

    \column{0.38\textwidth}
      \vspace{0.8em}
      \begin{itemize}
        \item $r_1(t) = 2\!\left(1 - \dfrac{t}{T}\right)$
        \vspace{0.5em}
        \item \good{$r_1 \geq 1$}: primeras $T/2$ iteraciones $\to$ exploración
        \vspace{0.3em}
        \item \accent{$r_1 < 1$}: últimas $T/2$ iteraciones $\to$ explotación
        \vspace{0.5em}
        \item División \textbf{50/50} automática con $a=2$
        \vspace{0.5em}
        \item[\warn{!}] El umbral real efectivo es $r_1^* = \pi/2 \approx 1{,}57$ porque $\mathbb{E}[|\phi|] = r_1 \cdot \frac{2}{\pi}$
      \end{itemize}
  \end{columns}
\end{frame}

% ----------------------------------------------------------------
% 3.3 Figura factor trigonométrico
% ----------------------------------------------------------------
\begin{frame}{$r_2$: dirección y signo del movimiento}
  \begin{columns}[T]
    \column{0.60\textwidth}
      \begin{tikzpicture}
      \begin{axis}[
        width=\textwidth, height=5.8cm,
        axis lines*=left,
        x axis line style={draw=none},
        y axis line style={draw=black, line width=0.6pt},
        xlabel={$r_2$}, ylabel={$r_1\cdot f(r_2)$},
        label style={font=\small\sffamily},
        tick label style={font=\scriptsize\sffamily},
        xmin=0, xmax=6.2832, ymin=-2.3, ymax=2.9,
        xtick={0, 1.5708, 3.1416, 4.7124, 6.2832},
        xticklabels={$0$,$\frac{\pi}{2}$,$\pi$,$\frac{3\pi}{2}$,$2\pi$},
        ytick={-2,-1,0,1,2},
        tick align=outside,
        grid=both,
        grid style={line width=0.4pt, draw=gray!30},
        legend style={font=\scriptsize\sffamily, at={(0.97,0.97)},
          anchor=north east, draw=none, fill=white,
          fill opacity=0.85, text opacity=1, row sep=2pt},
        legend cell align={left},
      ]
        \addplot[draw=black, line width=0.6pt, domain=0:6.2832, forget plot]{0};
        \addplot[fill=myyellow!40, opacity=0.6, draw=none, forget plot]
          coordinates {(0,-1)(6.2832,-1)(6.2832,1)(0,1)} \closedcycle;
        \addplot[very thick, myblue,           domain=0:6.2832, samples=200]{ 2*sin(deg(x))};
        \addlegendentry{$r_1\!=\!2$: $r_1\sin(r_2)$}
        \addplot[very thick, myblue,   dashed, domain=0:6.2832, samples=200]{ 2*cos(deg(x))};
        \addlegendentry{$r_1\!=\!2$: $r_1\cos(r_2)$}
        \addplot[very thick, myyellow,         domain=0:6.2832, samples=200]{0.5*sin(deg(x))};
        \addlegendentry{$r_1\!=\!0.5$: $r_1\sin(r_2)$}
        \addplot[very thick, myyellow, dashed, domain=0:6.2832, samples=200]{0.5*cos(deg(x))};
        \addlegendentry{$r_1\!=\!0.5$: $r_1\cos(r_2)$}
        \node[font=\footnotesize\bfseries, myyellow!80!black] at (axis cs:3.1416, 0.6)
          {zona explot.};
      \end{axis}
      \end{tikzpicture}

    \column{0.37\textwidth}
      \vspace{0.8em}
      \begin{itemize}
        \item $r_2 \in [0, 2\pi]$ determina el \textbf{signo} de $\phi$
        \vspace{0.3em}
        \item $f(r_2) > 0$: movimiento \textit{hacia} el destino
        \vspace{0.2em}
        \item $f(r_2) < 0$: movimiento \textit{alejándose} del destino
        \vspace{0.4em}
        \item \good{$r_1=2$ (azul)}: el factor sale de $[-1,1]$\\
              $\to$ exploración frecuente
        \vspace{0.3em}
        \item \accent{$r_1=0.5$ (naranja)}: siempre dentro de $[-1,1]$\\
              $\to$ explotación pura
        \vspace{0.4em}
        \item Sin y cos desfasados $\pi/2$ mejoran la \textbf{isotropía}
      \end{itemize}
  \end{columns}
\end{frame}

% ----------------------------------------------------------------
% 3.4 Figura geometría 4 regímenes
% ----------------------------------------------------------------
\begin{frame}{Cuatro regímenes de movimiento (1D)}
  \begin{columns}[T]
    \column{0.60\textwidth}
      \vspace{0.5em}
      \begin{tikzpicture}[>=Stealth, thick, font=\small]
        \coordinate (X) at (0,0);
        \coordinate (P) at (4,0);
        % Bandas
        \fill[myyellow!40, opacity=0.6] (-4,-0.38) rectangle (4, 0.38);
        \fill[mygreen!40,  opacity=0.6] ( 4,-0.38) rectangle (7, 0.38);
        \fill[mygreen!40,  opacity=0.6] (-4,-0.38) rectangle (-7,0.38);
        % Eje
        \draw[gray!50, thin] (-6.8,0) -- (6.8,0);
        % Cotas
        \draw[<->, gray!60, thin] (0,-0.75) -- (4,-0.75)
          node[midway, below, gray!80] {\footnotesize $|\Delta_j^t|$};
        \draw[dashed, gray!50, thin] (4,0) -- (4,-0.9);
        \draw[dashed, gray!50, thin] (0,0) -- (0,-0.9);
        % Flechas exploración
        \draw[->, thick, myblue]        (0,0) -- ( 5.8,0)
          node[above] {\footnotesize $\phi\!>\!1$};
        \draw[->, thick, mygreen!70!black] (0,0) -- (-5.5,0)
          node[above] {\footnotesize $\phi\!<\!-1$};
        % Flechas explotación
        \draw[->, very thick, myyellow!80!black] (0,0) -- ( 2.5,0)
          node[below] {\footnotesize $0\!<\!\phi\!<\!1$};
        \draw[->, very thick, myyellow]          (0,0) -- (-2.0,0)
          node[below] {\footnotesize $-1\!<\!\phi\!<\!0$};
        % Etiquetas zona
        \node[myyellow,         font=\footnotesize\bfseries] at ( 0,  0.7) {Explotación};
        \node[mygreen,          font=\footnotesize\bfseries] at ( 5.7,0.7) {Exploración};
        \node[mygreen!70!black, font=\footnotesize\bfseries] at (-5.7,0.7) {Exploración};
        % Puntos
        \filldraw[myred]  (P) circle (3pt);
        \node[above=4pt, myred]  at (P) {$P_j^t$};
        \filldraw[black]  (X) circle (3pt);
        \node[above=4pt, black]  at (X) {$X_j^t$};
      \end{tikzpicture}

    \column{0.37\textwidth}
      \vspace{0.5em}
      \small
      \begin{tabular}{c l}
        \toprule
        $\phi$ & Régimen \\
        \midrule
        $(0,\,1)$    & \accent{Explot. positiva} \\
        $(-1,\,0)$   & \accent{Explot. negativa} \\
        $>\!1$       & \good{Explor. positiva}   \\
        $<\!-1$      & \good{Explor. negativa}   \\
        \bottomrule
      \end{tabular}

      \vspace{0.6em}
      \begin{itemize}
        \item Simetría izq./der. $\Rightarrow$ \textbf{sin sesgo} de dirección
        \item Los 4 casos se dan por probabilidad en cada iteración
        \item La zona de explotación se \textit{estrecha} con $r_1$
      \end{itemize}
  \end{columns}
\end{frame}

% ----------------------------------------------------------------
% 3.5 Papel de r3 y r4 en el equilibrio
% ----------------------------------------------------------------
\begin{frame}{$r_3$ y $r_4$: modulación secundaria del equilibrio}
  \begin{columns}[T]
    \column{0.48\textwidth}
      \begin{block}{$r_3 \in [0,2]$: distancia efectiva al destino}
        \[|\Delta_j^t| = |r_3 P_j^t - X_j^t|\]
        \begin{itemize}
          \item $r_3 > 1$: desplaza el destino \textit{más lejos}, agranda $|\Delta_j^t|$\\
                $\to$ \good{efecto exploratorio secundario}
          \item $r_3 < 1$: acerca el destino a $X_j^t$, reduce $|\Delta_j^t|$\\
                $\to$ \accent{efecto explotador secundario}
          \item $r_3 = 1$: distancia real sin modificar
        \end{itemize}
        \textbf{Clave:} introduce diversidad \emph{incluso cuando} $r_1 < 1$
      \end{block}

    \column{0.48\textwidth}
      \begin{block}{$r_4 \in [0,1]$: selección sin/cos}
        \begin{itemize}
          \item Probabilidad $1/2$ para cada función
          \item $\sin$ y $\cos$ desfasadas $\pi/2$ rad
          \item Juntas cubren de forma más \textbf{uniforme} el círculo unitario
        \end{itemize}
        \vspace{0.3em}
        \begin{center}
          \begin{tikzpicture}[scale=0.9]
            \draw[gray!40] (0,0) circle (1cm);
            \draw[gray!40] (0,0) circle (1.6cm);
            % sin muestras
            \foreach \a in {0,30,...,330}{
              \draw[myblue!60, ->, thin]
                (0,0) -- ({1.6*sin(\a)}, {1.6*cos(\a)});
            }
            % cos muestras (desfasadas 90°)
            \foreach \a in {15,45,...,345}{
              \draw[myyellow!80, ->, thin]
                (0,0) -- ({1.6*sin(\a)}, {1.6*cos(\a)});
            }
            \node[font=\tiny, myblue]   at (1.9, 1.4)  {sin};
            \node[font=\tiny, myyellow] at (-1.9,1.4)  {cos};
            \filldraw[black] (0,0) circle (1.5pt);
            \node[font=\tiny] at (0,-0.25) {$X_j^t$};
          \end{tikzpicture}
        \end{center}
      \end{block}
  \end{columns}
\end{frame}

% ----------------------------------------------------------------
% 3.6 Resumen cuantitativo del equilibrio
% ----------------------------------------------------------------
\begin{frame}{Análisis cuantitativo del equilibrio}

  \begin{block}{Valor esperado del factor de movimiento}
    \[
      \mathbb{E}[|\phi|] = r_1 \cdot \mathbb{E}[|f(r_2)|] = r_1 \cdot \frac{2}{\pi}
    \]
    El factor $\tfrac{2}{\pi} \approx 0{,}637$ reduce el umbral \textit{efectivo}: la transición real ocurre en $r_1^* = \tfrac{\pi}{2} \approx 1{,}57$, es decir, en $t^* \approx 0{,}21\,T$
  \end{block}

  \vspace{0.3em}
  \begin{columns}[T]
    \column{0.52\textwidth}
      \begin{tikzpicture}
      \begin{axis}[
        width=\textwidth, height=4.5cm,
        axis lines*=left,
        xlabel={$t/T$}, ylabel={$\mathbb{E}[|\phi|]$},
        label style={font=\small\sffamily},
        tick label style={font=\scriptsize\sffamily},
        xmin=0, xmax=1, ymin=0, ymax=1.35,
        xtick={0,0.21,0.5,1},
        xticklabels={$0$, $0{,}21$, $0{,}5$, $1$},
        ytick={0,0.5,0.637,1},
        yticklabels={$0$,$0{,}5$,$\frac{2}{\pi}$,$1$},
        tick align=outside,
        grid=both,
        grid style={line width=0.3pt, draw=gray!25},
        legend style={font=\scriptsize\sffamily, at={(0.97,0.95)},
          anchor=north east, draw=none, fill=white,
          fill opacity=0.85, text opacity=1},
      ]
        \addplot[very thick, myblue, domain=0:1, samples=200]
          {2*(1-x)*2/pi};
        \addlegendentry{$\mathbb{E}[|\phi|]$}
        \addplot[dashed, myred, thick, domain=0:1, samples=2]{0.637};
        \addlegendentry{Umbral $\frac{2}{\pi}$}
        \addplot[fill=mygreen!30, opacity=0.55, draw=none]
          coordinates {(0,0.637)(0.21,0.637)(0.21,1.3)(0,1.3)} \closedcycle;
        \addplot[fill=myyellow!30, opacity=0.55, draw=none]
          coordinates {(0.21,0)(1,0)(1,0.637)(0.21,0.637)} \closedcycle;
        \node[font=\tiny\bfseries, mygreen]  at (axis cs:0.10, 1.1) {Explor.};
        \node[font=\tiny\bfseries, myyellow] at (axis cs:0.60, 0.3) {Explot.};
      \end{axis}
      \end{tikzpicture}

    \column{0.45\textwidth}
      \vspace{0.5em}
      \small
      \begin{tabular}{l l}
        \toprule
        Parámetro & Efecto en equilibrio \\
        \midrule
        $r_1 \to 0$ & Explotación total \\
        $r_1 = 1$   & Umbral nominal \\
        $r_1 = \pi/2$ & Umbral \textit{efectivo} real \\
        $r_1 = a=2$ & Exploración máxima \\
        \midrule
        $r_3 > 1$   & +Exploración \\
        $r_3 < 1$   & +Explotación \\
        $r_2$ unif. & Isotropía \\
        $r_4$ unif. & Mejor cobertura \\
        \bottomrule
      \end{tabular}

      \vspace{0.4em}
      \warn{La transición real ocurre aprox.\ en $t \approx 0{,}21\,T$, no en $t = 0{,}5\,T$}
  \end{columns}
\end{frame}

% ----------------------------------------------------------------
% 3.7 Diagrama síntesis del equilibrio
% ----------------------------------------------------------------
\begin{frame}[fragile]{Síntesis: flujo del equilibrio exploración/explotación}
  \begin{center}
  \begin{tikzpicture}[>=Stealth, font=\small,
    bloque/.style={rounded corners=4pt, draw, thick,
                   minimum width=2.8cm, minimum height=0.9cm,
                   align=center, fill=#1, text=white},
    flecha/.style={->, thick}
  ]
    % Nodos
    \node[bloque=myblue]    (r1)  at (0,0)      {$r_1(t)$ decrece\\linealmente};
    \node[bloque=mygreen]   (exp) at (-3.5,-2)  {Exploración\\$|\phi|>1$};
    \node[bloque=myyellow!80!black]  (expl) at ( 3.5,-2) {Explotación\\$|\phi|<1$};
    \node[bloque=myred!80]  (r3e) at (-3.5,-4)  {$r_3>1$\\amplía $|\Delta|$};
    \node[bloque=myblue!60] (r3x) at ( 3.5,-4)  {$r_3<1$\\reduce $|\Delta|$};
    \node[bloque=gray!60]   (r2)  at (0,-2)     {$r_2$: dirección\\$r_4$: sin/cos};
    \node[bloque=gray!50]   (iso) at (0,-4)     {Isotropía\\sin sesgo};

    % Flechas principales
    \draw[flecha, mygreen]         (r1) -- node[left,  black, font=\scriptsize]{$t < T/2$} (exp);
    \draw[flecha, myyellow!80!black] (r1) -- node[right, black, font=\scriptsize]{$t > T/2$} (expl);
    \draw[flecha, gray!70]         (r1) -- (r2);

    % Flechas secundarias r3
    \draw[flecha, myred!70]   (exp)  -- (r3e);
    \draw[flecha, myblue!60]  (expl) -- (r3x);
    \draw[flecha, gray!60]    (r2)   -- (iso);

    % Etiquetas de modulación
    \node[font=\scriptsize, myred!80]   at (-1.9,-3.1)  {+explor.};
    \node[font=\scriptsize, myblue!80]  at ( 1.9,-3.1)  {+explot.};

  \end{tikzpicture}
  \end{center}
\end{frame}

% ================================================================
\section{Pseudocódigo}
% ================================================================

\begin{frame}[fragile]{Pseudocódigo del SCA}
\begin{lstlisting}
(* Parámetros: N (población), T (iteraciones máx.), a=2 *)
Inicializar población X = {X_1,...,X_N} en [lb, ub]^d
Evaluar f(X_i) para cada i;   P <- argmin f(X_i)

t <- 1
Mientras t <= T hacer
    r1 <- a * (1 - t/T)                   (* Ec. decaimiento *)
    Para cada agente i = 1 hasta N hacer
        Para cada dimensión j = 1 hasta d hacer
            r2 <- U[0, 2*pi];  r3 <- U[0,2];  r4 <- U[0,1]
            Si r4 < 0.5 entonces
                X_i[j] <- X_i[j] + r1*sin(r2)*|r3*P[j] - X_i[j]|
            Si_no
                X_i[j] <- X_i[j] + r1*cos(r2)*|r3*P[j] - X_i[j]|
            Fin Si
            X_i[j] <- max(lb[j], min(ub[j], X_i[j]))
        Fin Para
        Si f(X_i) < f(P) entonces  P <- X_i  Fin Si
    Fin Para;   t <- t + 1
Fin Mientras
Retornar P
\end{lstlisting}
  \vspace{0.2em}
  \small Complejidad: $\mathcal{O}(N \cdot d)$ por iteración; $N \times T$ evaluaciones totales.
\end{frame}

% ================================================================
\section{Ventajas, limitaciones y aplicaciones}
% ================================================================

\begin{frame}{Ventajas e inconvenientes}
  \small
  \begin{columns}[T]
    \column{0.49\textwidth}
      \begin{exampleblock}{\good{Ventajas}}
        \begin{itemize}
          \item Implementación simple: solo $a$, $N$, $T$
          \item Balance explor./explot.\ \textbf{automático}
          \item Convergencia competitiva frente a PSO, GA, GWO
          \item Bajo coste: $\mathcal{O}(N \cdot d)$
          \item Fácilmente hibridable (PSO, DE, ABC)
        \end{itemize}
      \end{exampleblock}

    \column{0.49\textwidth}
      \begin{alertblock}{\warn{Inconvenientes}}
        \begin{itemize}
          \item Convergencia prematura en funciones multimodales
          \item Decaimiento lineal de $r_1$: no se adapta al paisaje
          \item Baja diversidad en iteraciones avanzadas
          \item Escasa memoria colectiva (solo $P_j^t$ compartido)
          \item Rendimiento pobre en espacios discretos sin adaptación
        \end{itemize}
      \end{alertblock}
  \end{columns}

  \vspace{0.5em}
  \begin{block}{Cuándo funciona bien}
    \begin{itemize}
      \item Problemas \textbf{continuos} de dimensión baja/media ($d \leq 100$)
      \item Funciones \textbf{unimodales} o con pocos óptimos locales
      \item Cuando la evaluación de $f$ es costosa (bajo overhead)
    \end{itemize}
  \end{block}
\end{frame}

\begin{frame}{Aplicaciones reales}
  \begin{columns}[T]
    \column{0.50\textwidth}
      \begin{block}{Ingeniería eléctrica}
        Flujo de potencia óptimo (OPF), ubicación de D-STATCOM, planificación hidrotérmica
      \end{block}
      \vspace{0.3em}
      \begin{block}{Aprendizaje automático}
        Selección de características (SCA binario), ajuste de SVR y redes neuronales
      \end{block}
      \vspace{0.3em}
      \begin{block}{Imagen médica}
        Umbralización multinivel en RM/TC, diagnóstico de cáncer de pulmón
      \end{block}

    \column{0.47\textwidth}
      \begin{block}{Ingeniería mecánica}
        Perfiles aerodinámicos, sistemas de engranajes, detección de daños estructurales
      \end{block}
      \vspace{0.3em}
      \begin{block}{Robótica y redes}
        Planificación de trayectorias (AUV), colocación de cámaras y sensores inalámbricos
      \end{block}
      \vspace{0.3em}
      \begin{block}{Variantes destacadas}
        \small SCA-PSO, OBSCA, SCA caótico, MO-SCA, SCA binario, EBSCA
      \end{block}
  \end{columns}
\end{frame}

% ================================================================
\section{Conclusiones}
% ================================================================

\begin{frame}{Conclusiones}
  \begin{enumerate}
    \item El SCA es una metaheurística \highlight{matemáticamente elegante}: seno y coseno proporcionan exploración e isotropía de forma natural.

    \vspace{0.4em}
    \item El parámetro $r_1$ es el \highlight{director del equilibrio}: su decaimiento lineal convierte la búsqueda en un proceso de \textbf{enfriamiento implícito}.

    \vspace{0.4em}
    \item El umbral efectivo de transición es $r_1^* = \pi/2$, lo que implica que la explotación predomina antes de $t = T/2$ en términos probabilísticos.

    \vspace{0.4em}
    \item $r_3$ actúa como modulador secundario: introduce diversidad \textbf{incluso en la fase de explotación}.

    \vspace{0.4em}
    \item Sus principales limitaciones (convergencia prematura, escasa memoria colectiva) han motivado una familia rica de variantes híbridas desde 2016.

    \vspace{0.4em}
    \item Como dicta el teorema \textit{No Free Lunch}: es competitivo en funciones continuas suaves, pero no universalmente superior.
  \end{enumerate}
\end{frame}

% ----------------------------------------------------------------
% CIERRE
% ----------------------------------------------------------------
\begin{frame}[plain]
  \begin{center}
    \vspace{2em}
    {\Large \textbf{Sine Cosine Algorithm}}\\[0.5em]
    {\large Jesús Muñoz Velasco}\\[0.3em]
    {\small Metaheurísticas · Curso 2025--2026}

    \vspace{2em}
    \begin{tikzpicture}
    \begin{axis}[
      width=9cm, height=3.5cm,
      axis lines=none,
      xmin=0, xmax=12.566, ymin=-2.3, ymax=2.3,
    ]
      \addplot[very thick, myblue,  domain=0:12.566, samples=300]{2*sin(deg(x))};
      \addplot[very thick, myred,   domain=0:12.566, samples=300]{1.5*cos(deg(x))};
      \addplot[very thick, mygreen, domain=0:12.566, samples=300]{sin(deg(2*x))};
    \end{axis}
    \end{tikzpicture}

    \vspace{1em}
    {\small \textit{``SCA: A Sine Cosine Algorithm for solving optimization problems''}}\\
    {\small S.~Mirjalili, Knowledge-Based Systems, 2016}
  \end{center}
\end{frame}

\end{document}